\section{Preliminaries}
 
For a function $f\colon X\to\R$ on a finite set 
$X,$ we let $\|f\|_{\infty}=\max_{x\in X}|f(x)|$.
The total degree of a multivariate real polynomial $p\colon\R^n\to\R$
is denoted $\deg p$. We use the terms \emph{degree} and \emph{total
degree} interchangeably in this paper. For a function $f\colon X\to\R$
on a finite set $X\subset\R^n$, the \emph{$\epsilon$-approximate
degree} $\degeps(f)$ of $f$ is defined as the least degree of a
real polynomial $p$ with $\|f-p\|_{\infty}\leq\epsilon$. Throughout
this paper, we will work with the $\epsilon$-approximate degree for
a small constant $\epsilon>0$. For Boolean functions $f\colon X\to\zoo$,
the choice of constant $0<\epsilon<1/2$ affects the quantity $\degeps(f)$
by at most a constant factor:
\begin{align}
c\adeg(f)\leq\degeps(f)\leq C\adeg(f)\,,\label{eq:error-boosting}
\end{align}
where $c=c(\epsilon)$ and $C=C(\epsilon)$ are positive constants.
By convention, one studies $\epsilon=1/3$ as the canonical case and
reserves for it the special symbol $\tildeg(f)=\adeg(f)$. A dual
characterization~\cite{sherstov07quantum,sherstov11quantum-sdpt}
of the approximate degree is as follows.

\begin{fact}
\label{fact:dual}
Let $f\colon X\to\R$ be given, for a finite set
$X\subset\R^n$. Then $\degeps(f)\geq d$ if and only if there
exists a function $\psi\colon X\to\R$ such that
\begin{align*}
& \sum_{x\in X}|\psi(x)|=1\,,\\
& \sum_{x\in X}\psi(x)f(x)>\epsilon\,,
\intertext{and}
&\sum_{x\in X}\psi(x)p(x)=0
\end{align*}
for every polynomial $p$ of degree less than $d$.
\end{fact}

We adopt the usual definitions of the Boolean functions 
$\AND_n,\OR_n\colon\zoon\to\zoo$.
Their approximate degree was determined by Nisan 
and Szegedy~\cite{nisan-szegedy94degree}.

\begin{theorem}[Nisan and Szegedy]
\label{thm:nisan-szegedy}The functions $\AND_n$ and $\OR_n$
obey
\begin{align*}
\adeg(\AND_n)=\adeg(\OR_n)=\Theta(\sqrt{n})\,.
\end{align*}
\end{theorem}

By combining the above two theorems, Gavinsky and the author~\cite[Thm.~5.1]{GS09npconp}
obtained the following result, which plays a key role in this paper.

\begin{theorem}[Gavinsky and Sherstov]
\label{thm:dual}Fix any constant $0<\epsilon<1$. Then there exists
a constant $\delta=\delta(\epsilon)>0$ and a real function 
$\psi\colon\zoon\to\R$ such that
\begin{align}
& \sum_{x\in\zoon}|\psi(x)|=1\,,\label{eq:psi-bounded}\\
& \psi(0,0,\dots,0)<-\frac{1-\epsilon}2\,,\label{eq:psi-metric}
\intertext{and}
&\sum_{x\in\zoon}\psi(x)p(x)  =0\label{eq:psi-orthog}
\end{align}
for every polynomial $p$ of degree less than $\delta\sqrt{n}$.
\end{theorem}

For the sake of completeness, we include the proof.

\begin{proof}[Proof of \expref{Theorem}{thm:dual} 
(adapted from~\cite{GS09npconp})]
Recall from \expref{Theorem}{thm:nisan-szegedy} that $\adeg(\OR_n)=\Omega(\sqrt{n})$.
Therefore, \eqref{eq:error-boosting} shows that 
$\deg_{\frac{1-\epsilon}2}(\OR_n)\geq\delta\sqrt{n}$
for a sufficiently small constant $\delta=\delta(\epsilon)>0$. Now
the dual characterization of the approximate degree (\expref{Fact}{fact:dual})
provides a function $\psi\colon\zoon\to\R$ that obeys \eqref{eq:psi-bounded},
\eqref{eq:psi-orthog}, and
\begin{align}
\sum_{x\in\zoon}\psi(x)\OR_n(x)>\frac{1-\epsilon}2\,.\label{eq:correl-with-OR}
\end{align}
It remains to verify \eqref{eq:psi-metric}:
\begin{align*}
\psi(0,0,\dots,0) & =\sum_{x\in\zoon}\psi(x)(1-\OR_n(x))\\
& =-\sum_{x\in\zoon}\psi(x)\OR_n(x) & & \text{by \eqref{eq:psi-orthog}}\\
& <-\frac{1-\epsilon}2 & & \text{by
\eqref{eq:correl-with-OR}.}\qedhere
\end{align*}
\end{proof}

For probability distributions $\mu$ and $\lambda$ on finite sets
$X$ and $Y$, respectively, we let $\mu\times\lambda$ denote the
probability distribution on $X\times Y$ given by 
$(\mu\times\lambda)(x,y)=\mu(x)\lambda(y)$.
The \emph{support }of a probability distribution $\mu$ is defined
to be $\supp\mu=\{x:\mu(x)>0\}$.
 
 
\section{Main Result}
 
We are now in a position to prove our main result.

%
%
\begin{theorem}
The Boolean function $f(x)=\bigwedge_{i=1}^m\bigvee_{j=1}^nx_{ij}$
obeys
\begin{align}
\adeg(f)=\Omega(\sqrt{mn})\,.\label{eq:f-adeg-main-equation}
\end{align}
%
\end{theorem}

\begin{proof}
Let $\epsilon$ be an absolute constant to be named later, $0<\epsilon<1$.
Then by \expref{Theorem}{thm:dual}, there exists a constant $\delta=\delta(\epsilon)>0$
and a function $\psi\colon\zoon\to\R$ that
obeys~\eqref{eq:psi-bounded}--\eqref{eq:psi-orthog}.
Let $\mu$ be the probability distribution on $\zoon$ given by $\mu(x)=|\psi(x)|$.
Let $\mu_0$ and $\mu_1$ be the probability distributions induced
by $\mu$ on the sets $\{x:\psi(x)<0\}$ and $\{x:\psi(x)>0\},$ respectively.
Since $\sum_{x\in\zoon}\psi(x)=0,$ the sets $\{x:\psi(x)<0\}$ and
$\{x:\psi(x)>0\}$ are weighted equally by $\mu$. As a consequence,
\begin{align}
\mu & =\frac12\mu_1+\frac12\mu_0\,,\label{eq:mu-mu0-mu1}\\
\psi & =\frac12\mu_1-\frac12\mu_0\,.\label{eq:psi-mu0-mu1}
\end{align}
 
 
Consider the linear operator $L$ that maps functions $\phi\colon(\zoon)^m\to\R$
to functions $L\phi\colon\zoom\to\R$ according to
\begin{align*}
(L\phi)(z)=\Exp_{x_1\sim\mu_{z_1}}\cdots\Exp_{x_m\sim\mu_{z_m}}\phi(x_1,\dots,x_m)\,.
\end{align*}
Fix a real polynomial $p$ with
\begin{align}
\|f-p\|_{\infty}\leq\epsilon\,.\label{eq:F-p}
\end{align}
 
\begin{claim} \label{cla:Lf-AND}
\quad $\|\AND_m-Lf\|_{\infty}<\epsilon$.
\end{claim}
 
\begin{claim}
\label{cla:deg}
\quad $\deg p\geq\delta\sqrt{n}\;\deg Lp$.
\end{claim}

Before settling the claims, we finish the proof of the theorem. The
linearity of $L$ yields
\begin{align*}
\|\AND_m-Lp\|_{\infty} & \leq\underbrace{\|\AND_m-Lf\|_{\infty}}_{<\epsilon}
+\underbrace{\|L(f-p)\|_{\infty}}_{\leq\epsilon}<2\epsilon\,,
\end{align*}
where we have used~\eqref{eq:F-p} and \expref{Claim}{cla:Lf-AND} in
bounding the marked quantities. For $\epsilon=1/6,$ we arrive at
$\|\AND_m-Lp\|_{\infty}\leq1/3$ and therefore $\deg Lp=\Omega(\sqrt{m})$
by \expref{Theorem}{thm:nisan-szegedy}. Now \expref{Claim}{cla:deg} implies
that $\deg p=\Omega(\sqrt{mn})$.
\end{proof}
 
\begin{proof}[Proof of \expref{Claim}{cla:Lf-AND}]
By~\eqref{eq:psi-metric}, we have $\psi(x)>0$ only when $\OR_n(x)=1$.
Hence $\supp\mu_1\subseteq\OR_n^{-1}(1)$ and
\begin{align*}
(Lf)(1,1,\dots,1)
   =\Exp_{\mu_1\times\cdots\times\mu_1}[f]
   =\prod_{i=1}^m\Exp_{\mu_1}[\OR_n]=1\,.
\end{align*}
 
 
It remains to prove that $|(Lf)(z)|<\epsilon$ for every $z\ne(1,1,\dots,1)$.
We have
\begin{align*}
(Lf)(z)
  =\Exp_{\mu_{z_1}\times\cdots\times\mu_{z_m}}[f]
  =\prod_{i=1}^m\Exp_{\mu_{z_i}}[\OR_n]
  =\prod_{i=1}^m(1-\mu_{z_i}(0,0,\dots,0))\,,
\end{align*}
whence
\begin{align}
0\leq(Lf)(z)\leq1-\mu_0(0,0,\dots,0)\,.\label{eq:Lf-0-1-mu}
\end{align}
We know from~\eqref{eq:psi-metric} that $\psi(0,0,\dots,0)<-(1-\epsilon)/2,$
which means in particular that $(0,0,\dots,0)\in\supp\mu_0$. Therefore
\begin{align*}
\mu_0(0,0,\dots,0)
  =2\mu(0,0,\dots,0)
  =2|\psi(0,0,\dots,0)|
  >1-\epsilon\,,
\end{align*}
where the first step uses \eqref{eq:mu-mu0-mu1}. By~\eqref{eq:Lf-0-1-mu},
we conclude that $0\leq(Lf)(z)<\epsilon$.
\end{proof}
 
\begin{proof}[Proof of \expref{Claim}{cla:deg}]
By the linearity of $L$, it suffices to consider factored polynomials
$p$ of the form $p(x)=\prod_{i=1}^mp_i(x_{i,1},x_{i,2},\dots,x_{i,n})$.
In this case we have the convenient formula
\begin{align*}
(Lp)(z)=\prod_{i=1}^m\Exp_{\mu_{z_i}}[p_i]\,.
\end{align*}
By \eqref{eq:psi-orthog} and \eqref{eq:psi-mu0-mu1}, polynomials
$p_i$ of degree less than $\delta\sqrt{n}$ obey 
$\Exp_{\mu_0}[p_i]=\Exp_{\mu_1}[p_i]$
and therefore do not contribute to the degree of $Lp$. As a result,
\[
\deg Lp\leq|\{i:\deg p_i\geq\delta\sqrt{n}\}|
   \leq\frac{\deg p}{\delta\sqrt{n}}\,.\qedhere
\]
\end{proof}
